\begin{table}[h]
\centering
\begin{tabular}{ll}
\hline\hline
\noalign{\smallskip}
{\bf V0 from cascade selection criterion} & Value                      \\
\noalign{\smallskip}
\hline
\noalign{\smallskip}
DCA (baryon to PV)                         & $>0.03$ cm                           \\  
DCA (meson to PV)                          & $>0.04$ cm                           \\ 
DCA (h$^-$ to h$^+$)                       & $<1.5$ standard deviations                  \\
$\Lambda$ mass ($m_{\mathrm{V0}}$)         & $1.108<m_{\mathrm{V0}}<$ 1.124 \GeVmass     \\
Fiducial volume (R$_{2D}$)                 & $>$ 1.2 (1.1) cm                    \\
V0 pointing angle                          & $\cos(\theta_{\textrm{V0}})>$ 0.97                     \\

\noalign{\smallskip}
\hline\hline
\noalign{\smallskip}
{\bf Cascade finding criterion} & Value \\
\hline
\noalign{\smallskip} 
DCA (bachelor to PV)       & $>0.04$ cm\\
DCA (V0 to PV)                            & $>0.06$ cm\\
DCA ($\pi^{\pm}$ (K$^{\pm}$) to V0)       & $< 1.3$ cm\\    
Fiducial volume (R$_{2D}$)                & $>$ 0.6 (0.5) cm\\
Cascade pointing angle                    & $\cos(\theta_{\textrm{casc}})>$ 0.97\\
Proper Lifetime                  & $<$ 3$c\tau$ \\
Competing Cascade Rejection Window (\Oms\ only)                    & $\pm$8 \MeVmass\\
\noalign{\smallskip}                     
\hline\hline
\end{tabular}
~\newline
\caption{Selection criteria for V0 ($\Lambda$) from cascades, and cascades (\Xis\ and \Oms) presented in
this work. If a criterion for \Xis\ and \Oms\ finding differs, the criterion for \Oms\ hypothesis
is given in parentheses.
DCA stands for ``distance of closest approach" and PV for ``primary event vertex". The pointing angle
$\theta$ is the angle between the momentum vector of the reconstructed V0 or cascade and the line
segment bound by the decay and primary vertices  and R$_{2D}$ denotes the transverse distance from the 
detector center. 
The cascade track curvature is neglected, and $\tau$ refers to the average
lifetime for the two different cascade species. 
\label{tab:CascSels}
}
\end{table}
